\chapter{Conclusion}

This thesis studied Legal Judgment Prediction for Indian court cases as a structured prediction problem rather than a plain document-classification task. Credible legal prediction requires more than strong accuracy numbers: it requires a representation that preserves legally meaningful structure while actively resisting leakage from decision-bearing text and identity-based shortcuts.

The thesis introduced an end-to-end pipeline that converts raw Indian judgments into leakage-audited, graph-ready records through entity extraction, rhetorical-role segmentation, outcome labelling, multi-hearing consolidation, and reasoning-focused graph construction. On top of this representation, a Heterogeneous Graph Transformer modelled both the internal structure of each case and the shared legal links across cases via statutes, provisions, and precedents.

Across more than 73{,}000 cases spanning five legal domains, the framework achieved $0.8020$ accuracy, $0.7959$ macro-F1, and $0.8831$ ROC-AUC under 5-fold cross-validation on the best \texttt{party\_args\_lr\_decay} configuration, consistently outperforming text-only references. A probing classifier isolates a $+3.97$-point accuracy gain attributable to graph message passing at fixed classifier capacity --- evidence that gains come from cross-case legal structure, not surface text patterns.

A four-level interpretability pipeline --- representation probing, component importance, error decomposition, and case-level \path{Graph_Analyser} (heterogeneous PGExplainer + FAISS retrieval + LLM verbalisation) --- produces grounded, legally named explanations for individual predictions and exposes a clean structural signature of the model's confident failures: when wrong, the retrieved representation-space neighbours almost always reinforce the wrong label. A zero-shot evaluation on 11{,}769 food-safety judgments reaches $0.608$ accuracy and $0.675$ ROC-AUC against a $0.543$ majority baseline without any fine-tuning, showing that the cross-bucket HGT learns a partially domain-general legal-reasoning prior rather than memorising training-domain statutory hubs.

The main contribution is methodological: legal outcome prediction over Indian judgments becomes more credible when built around leakage-safe data construction, heterogeneous graph design, explicit constraints on what may be shared across cases, per-case explainability tied to the same representation the model predicts from, and an out-of-domain evaluation that tests whether the learned signal is about legal reasoning or corpus memorisation. While improvements are measured rather than dramatic, they support a clear conclusion: graph-based legal prediction is most useful when it is not only accurate, but structurally faithful, auditable at the case level, and experimentally honest about its transfer behaviour.

Future work can build on this foundation by tightening leakage audits, extending the outcome space beyond binary decisions, designing models that adapt to domain-specific difficulty and long factual narratives, closing the in-domain/out-of-domain gap through fine-tuning or domain-adaptive heads, and pairing the case-level explainer with a legal-expert validation loop to convert it from an internal diagnostic into an auditable review tool.
